%% ============================================================================
%% ch7_risks.tex -- Chapter 7: ניתוח סיכונים ותוכנית הפחתה
%% budget: body pages 35-36 (2 pages). R = P x I, critical if R >= 12 (ISO 31000)
%% v2 (red-team point 8): +RSK-ALG-07 (P95 miss), +RSK-OPS-08 (ethics/recruitment),
%% evidence for every residual probability, owners and early-warning indicators.
%% ============================================================================
\chapter{ניתוח סיכונים ותוכנית הפחתה}
\label{ch:risk}

\begin{managerbox}
סיכון שזוהה מראש הוא משימה בלוח הזמנים. סיכון שלא זוהה הוא משבר. הפרק מזהה שמונה סיכונים בחמישה תחומים, מכמת אותם, וקובע לכל סיכון קריטי פעולת מנע, תוכנית מגירה, בעלים ומדד התרעה, לפי תקן \en{ISO 31000} \cite{iso31000}.
\end{managerbox}

לכל תרחיש כשל מוערכות ההסתברות $P$ והחומרה $I$, בסולם \num{1} עד \num{5}, וציון הסיכון הוא $R = P \times I$ (משוואה~\ref{eq:risk}). סיכון שציונו \num{$R \ge 12$} מוגדר קריטי. איור~\ref{fig:heat} ממקם את הסיכונים, וטבלה~\ref{tab:risks} מפרטת אותם.
\begin{equation}
R = P \times I, \qquad P, I \in \{1,\dots,5\}
\label{eq:risk}
\end{equation}

\begin{reportfigure}{מטריצת סיכונים \num{5×5}: הסתברות מול חומרה, באדום האזור הקריטי (\num{$R \ge 12$})}{fig:heat}
\begin{tikzpicture}[x=0.95cm,y=0.6cm]
  \foreach \p in {1,...,5} \foreach \i in {1,...,5} {
    \pgfmathtruncatemacro{\r}{\p*\i}
    \ifnum\r>11 \def\col{brandwarn!70}\else\ifnum\r>5 \def\col{brandaccent!45}\else\def\col{brandgreen!30}\fi\fi
    \fill[\col] (\i-1,\p-1) rectangle (\i,\p);
    \draw[white,thick] (\i-1,\p-1) rectangle (\i,\p);
    \node[font=\tiny,text=black!45,anchor=north east] at (\i-0.04,\p-0.02) {\r};}
  \foreach \k in {1,...,5} {\node[font=\scriptsize] at (\k-0.5,-0.35) {\k}; \node[font=\scriptsize] at (-0.3,\k-0.5) {\k};}
  \node[font=\scriptsize\bfseries] at (2.5,-0.95) {Impact $I$};
  \node[font=\scriptsize\bfseries,rotate=90] at (-0.95,2.5) {Probability $P$};
  \foreach \i/\p/\n/\dy in {4/4/ALG-01/0, 3/4/GEN-02/0.14, 3/4/OPS-08/-0.2, 4/3/DAT-03/0.14, 4/3/ALG-07/-0.2, 5/3/BIZ-04/0, 5/2/SEC-05/0, 3/3/OPS-06/0}
    \node[font=\tiny\bfseries,fill=white,rounded corners=1pt,inner sep=0.8pt] at (\i-0.5,\p-0.55+\dy) {\n};
  \fill[brandwarn!70] (5.6,4.2) rectangle (6.0,4.6); \node[font=\tiny,anchor=west] at (6.05,4.4) {critical, $R\ge12$};
  \fill[brandaccent!45] (5.6,3.5) rectangle (6.0,3.9); \node[font=\tiny,anchor=west] at (6.05,3.7) {medium, $6\le R\le 11$};
  \fill[brandgreen!30] (5.6,2.8) rectangle (6.0,3.2); \node[font=\tiny,anchor=west] at (6.05,3.0) {low, $R\le 5$};
\end{tikzpicture}
\end{reportfigure}

מהאיור עולה שששה מתוך שמונת הסיכונים קריטיים, וכולם בחומרה \num{3} ומעלה. \en{RSK-SEC-05} אינו קריטי לפי הציון, אך חומרתו מרבית, ולכן הוא מטופל בתכן עצמו (\en{REQ-SEC-11}). \en{RSK-OPS-06} כולל גם את האפשרות שלא יימצא שותף אמיתי, כי השותף בדוח מדומה.

\begin{reporttable}[\footnotesize]{זיהוי הסיכונים וכימותם}{tab:risks}{P{1.85cm} R P{1.3cm} Q{0.4cm} Q{0.4cm} Q{0.55cm}}
\headrow \hc{מזהה} & \hc{תרחיש הכשל} & \hc{תחום} & \hc{$P$} & \hc{$I$} & \hc{$R$} \\ \midrule
\en{RSK-ALG-01} & תאורה קיצונית, עדשות מגע או עיניים עצומות מונעות כיול & אלגוריתמי & \num{4} & \num{4} & \textbf{\num{16}} \\
\en{RSK-GEN-02} & הזיות של מודל ההדמיה בידיים ובציפורניים & גנרטיבי & \num{4} & \num{3} & \textbf{\num{12}} \\
\en{RSK-DAT-03} & הגבלות גישה וקצב בממשקי הרשתות החברתיות & נתונים & \num{3} & \num{4} & \textbf{\num{12}} \\
\en{RSK-BIZ-04} & עלות \en{GPU} גבוהה מהתקרה או ירידה בהכנסות הפרסום & עסקי & \num{3} & \num{5} & \textbf{\num{15}} \\
\en{RSK-SEC-05} & דליפה של מידע ביומטרי & אבטחה & \num{2} & \num{5} & \num{10} \\
\en{RSK-OPS-06} & קטלוג דליל, התחלה קרה או היעדר שותף אמיתי & תפעולי & \num{3} & \num{3} & \num{9} \\
\en{RSK-ALG-07} & האחוזון ה-\num{95} של השגיאה לא יעמוד ב-\num{4.5} & אלגוריתמי & \num{3} & \num{4} & \textbf{\num{12}} \\
\en{RSK-OPS-08} & עיכוב באישור האתיקה או בגיוס המשתתפות & תפעולי & \num{4} & \num{3} & \textbf{\num{12}} \\
\end{reporttable}

לכל סיכון קריטי נקבעו פעולת מנע (\en{Mitigation}) ותוכנית מגירה (\en{Contingency Plan}). הציון השיורי מבוסס על ראיה, ולא על הצהרה: לכל ירידה בהסתברות מצוינת הבדיקה שתאשר אותה (טבלה~\ref{tab:mitig}).

\begin{reporttable*}[\scriptsize]{פעולות מנע, תוכניות מגירה, ציון שיורי והראיה לו, לסיכונים הקריטיים}{tab:mitig}{P{1.05cm} R R P{2.75cm}}
\headrow \hc{מזהה} & \hc{פעולת מנע} & \hc{תוכנית מגירה} & \hc{$R_{\text{residual}}$ והראיה לו} \\ \midrule
\en{ALG-01} & רישום בהבזק מסך; זיהוי עדשות ועיניים עצומות & למשתמשות שזוהו בלבד: כיול עם דף לבן, העוגן נשמר לשאר & \num{6 (2×3)}; שיעור הכשל בניסוי הבסיס \\
\en{GEN-02} & מסיכת יד וציפורניים כהתניה; סינון \en{LPIPS} & הדמיית פלג גוף עליון, ציפורניים כשכבה דו-ממדית & \num{6 (2×3)}; שיעור פסילה בשחזור, שבוע \num{5} \\
\en{DAT-03} & ממשקים רשמיים ותוכן ציבורי; מטמון \num{6} שעות & ספריית טרנדים עונתית ידנית & \num{4 (2×2)}; יומן חסימות, שבועות \num{3--4} \\
\en{BIZ-04} & זיקוק ל-\num{8} צעדים; תקרה \num{\$0.023}; ארבעה ערוצי הכנסה & רזולוציה נמוכה, או הדמיה בפרימיום בלבד & \num{6 (2×3)}; עלות מדודה, \en{TSK-4.3} \\
\en{ALG-07} & ניתוח הזנב לפי קבוצת גוון; עוגן שיניים כשלישי & דיווח שקוף; יעד הזנב עובר להמשך המחקר & \num{8 (2×4)}; רווח סמך באבן דרך א' \\
\en{OPS-08} & הגשה לאתיקה מיד עם האישור; מאגר משתתפות משני קמפוסים & ניסוי הבסיס על צוות ומתנדבות בהסכמה & \num{6 (2×3)}; אישור בכתב לפני שבוע \num{1} \\
\end{reporttable*}

הטבלה מראה שאף ציון שיורי אינו נשאר באזור הקריטי: ארבעה יורדים לאזור הבינוני, ואחד, \en{DAT-03}, לאזור הנמוך. תוכנית המגירה של \en{ALG-01} אינה מוותרת על הרעיון המרכזי: הדף הלבן מופעל רק אצל המשתמשות שאצלן העין אינה זמינה, וגם אז הוא משלים את העוגן ולא מחליף אותו. לכל סיכון יש בעלים, הסטודנטית שאחראית על המשימה שבה הוא נבדק לפי טבלה~\ref{tab:raci}, ומדד ההתרעה שבעמודה האחרונה נבחן בכל מפגש מנחה, לצד מדדי הבקרה בטבלה~\ref{tab:kpi}.
